\section{Tokenized-to-Parametric Transformations}
\label{sec:tok2param}

The transformation of experience from a Tokenized carrier to a Parametric carrier comprises two pathways: Evaluator Internalization (Section~\ref{sec:evaluator-internalization}) fixes interaction experience into an evaluation capability, and Policy Internalization (Section~\ref{sec:policy-internalization}) fixes interaction experience into a decision capability. Both pathways share one kind of source---the interaction trajectories produced by an agent's sequential decision-making.

\subsection{Evaluator Internalization: Tokenized-to-Evaluator Transformation}
\label{sec:evaluator-internalization}

Evaluator Internalization internalizes tokenized agent experience into Evaluator parameters, fixing the evaluation capability in the weights. Its products span reward models (outcome and process), verifiers, critics, and failure detectors. The inclusion premise is that the Evaluator parameters are substantively updated by experience data; an LLM-as-a-judge triggered purely by a prompt, with parameters unchanged, is excluded.

\subsubsection{Outcome-supervised Evaluator Internalization}

Outcome-supervised internalization binds the supervision signal to a complete response, an entire trajectory, or a whole episode, and maps it to a single evaluative label. By output form, this subsection splits into discriminative and generative outcome evaluators.

\paragraph{Discriminative Outcome Evaluators.} These output a scalar verdict, split by signal type into binary outcome judgment and graded outcome scoring.

\textbf{Binary outcome judgment} maps a complete episode to a binary verdict such as success/failure or violation/compliance. Du et al.~\cite{Du23} cast success detection as visual question answering, training Flamingo to make a binary judgment from clip-level success labels. I-FailSense~\cite{Gri25} extends this to cross-task failure detection, with positives from the original instruction and negatives from semantically mismatched instructions. WebRL~\cite{Qi24} trains an outcome reward model within a self-evolving curriculum RL, judging whether a web trajectory completes its task from the instruction, action history, and page state. Wen et al.~\cite{Wen25b} move the target from task success to policy violation, judging violations against four rule classes one by one.

\textbf{Graded outcome scoring} gives a finer signal beyond binary---a graded score, a continuous probability, or a relative preference. RoboReward~\cite{Lee26b} performs counterfactual relabeling and negative clipping on real rollouts, building a 1--5 rubric that regresses onto discretized progress. SAFE~\cite{Gu25b} takes a VLA policy's internal latent as input and learns a time-varying failure probability over a rollout prefix, triggering abort by a conformal threshold. VLP~\cite{Liu25w} instead learns a relative preference, building comparison labels from expert/medium/random trajectories and learning a cross-modal preference under Bradley--Terry.

\paragraph{Generative Outcome Evaluators.} These output natural-language text in addition to a verdict, split by the role of the text into verdict justification, failure diagnosis, and corrective critique.

\textbf{Verdict justification} generates the reasoning that supports a verdict alongside it. GenRM~\cite{Zha24o} casts reward modeling as next-token prediction, with a GenRM-Direct that emits Yes/No directly and a GenRM-CoT that generates a CoT rationale before judging. S2J~\cite{Sun25l} has the judge solve the problem itself before judging, jointly optimizing the solving and judging rewards. J4R~\cite{Xu25i} has the judge keep a consistent verdict across equivalent inputs such as a response-order swap, suppressing position bias. AgentV-RL~\cite{Zha26ac} builds the verifier as a multi-turn tool-augmented agent, rewarding the consistency between verdict and ground truth. CriticGPT-RM~\cite{Liu24g} generates pairwise analyses and preferences automatically from operation videos at different stages, with a downstream independent reward model trained under Bradley--Terry.

\textbf{Failure diagnosis} generates an explanation or localization of a failure alongside the verdict. AHA~\cite{Dua24} injects multiple failure modes into successful demonstrations via FailGen, training the model to output both a yes/no verdict and a free-text failure rationale. ARMOR~\cite{Qi26} jointly trains a detection head and a reasoning decoder, generating detection and reasoning over several turns and selecting the result by confidence. ExeVRM~\cite{Son26d} trains on execution videos, outputting a video-level success judgment and a failure localization.

\textbf{Corrective critique} generates corrective feedback for a downstream generator or actor to consume. CTRL~\cite{Xie25c} optimizes the critique with GRPO under a unit-test-pass reward, making the generator more likely to produce code that passes the tests. Co-Evolving Critics~\cite{Li26l} has the critic give a diagnosis and the actor rewrite the trajectory accordingly, with the critic's reward including a saturation-aware gain to keep it from going stale as the policy evolves. RL4F~\cite{Aky23} trains a critique generator under a reward defined by the task metric of the refined output against ground truth, an early form of critique-driven repair.

\subsubsection{Process-supervised Evaluator Internalization}

Process-supervised internalization assigns the supervision signal to intermediate steps or action prefixes, writing local decision quality into the Evaluator parameters. By output form, this subsection splits into discriminative and generative process evaluators. The motivation comes from an empirical observation: the failure of a long-horizon agent is usually the accumulation of several local errors, and a single outcome label cannot separate those errors from a trajectory that fails as a whole.

\paragraph{Discriminative Process Evaluators.} These output a scalar score for each intermediate step or action prefix, split by signal semantics into step correctness scoring, step advantage and progress estimation, and step-action compatibility.

\textbf{Step correctness scoring} maps each step to a correct/incorrect label or a finer error-type label. Let's Verify Step by Step~\cite{Lig23} trains a PRM on PRM800K with human step-level labels, predicting positive/negative/neutral for each step. OmegaPRM~\cite{Luo24} uses MCTS and binary search to locate the first error automatically, replacing human annotation with a Monte Carlo estimate and producing soft process labels. PathFinder-PRM~\cite{Pal25} argues that reward estimation should be conditioned on error type---the model first decides explicitly whether the step made a math error or a consistency error, then predicts a step reward conditioned on that error type, rather than jumping from step text to a score. GUI-Shepherd~\cite{Che25h} brings step-correctness labels to the GUI domain, with humans labeling whether a state--action pair is correct and GPT-4o adding a CoT rationale as a training augmentation, outputting a scalar correctness score at inference. GAIA~\cite{Wan26p} organizes training as a data flywheel---a first round trains the critic by comparing correct and incorrect actions against ground truth, and a second round feeds the critic's hard cases back for a second round, continually absorbing the error distribution of a real GUI agent.

\textbf{Step advantage and progress estimation} gives a finer signal beyond binary: estimating one step's marginal contribution to the future success probability, or a continuous task-progress value. Rewarding Progress~\cite{Set24} unrolls a Monte Carlo rollout from each prefix to estimate the change in future success probability, letting the Evaluator learn the step advantage directly rather than step correctness. PQM~\cite{Li24m} formalizes the problem as Q-value ranking---replacing independent BCE with a Plackett--Luce ranking loss to force the estimated success probability of a correct prefix above that of an incorrect one. BiRM~\cite{Che25s} jointly trains a backward-looking reward head and a forward-looking value head on the same backbone, outputting both the local correctness of a step and the probability that a partial trajectory still leads to the correct answer. AgentPRM~\cite{Xi25b} extends this idea from reasoning to multi-step agent decisions---using TD and GAE to estimate the Q-value and advantage on state--action pairs from rollouts automatically, so that the intermediate steps of an agent trajectory can also receive a process reward. ProgRM~\cite{Zha25z} uses LCS recipe matching in GUI trajectories to locate key steps automatically and assign a normalized progress label, predicting task progress from the action history and current observation. Agentic-Q~\cite{Wan26s} lets the terminal success signal propagate back along the trajectory to train a step-wise success-probability estimator. SWE-Shepherd~\cite{Dih26} faces a particular difficulty in the code domain---a code-modification trajectory lacks a natural process-supervision signal---and so builds a weak process reward from heuristics such as relevant-file access, target-file modification, and test-result change, accumulated with discounting into a step-wise target to train a repo-level PRM.

\textbf{Step-action compatibility} introduces another signal beyond binary correctness: comparing the relative quality of several candidate actions at the same decision point. V-Droid~\cite{Dai25} centers on Pairwise Process Preference---rather than scoring a single action in isolation, it forces the correct action at a step to score above several incorrect ones, learning from the relative order rather than the absolute value. IntentScore~\cite{Che26s} extends the condition of compatibility from state to planning intent---the action encoder receives the planning intent explicitly as input and learns a compatibility score for state--action--intent triples, letting the evaluator distinguish ``the action is correct in itself'' from ``the action is correct under the current intent.''

\paragraph{Generative Process Evaluators.} These produce natural-language text---critique, error diagnosis, or corrective suggestion---in addition to a scalar verdict, split by the downstream function of the text into step-level critique, verdict with structured reasoning, and pre-operative corrective feedback.

\textbf{Step-level critique} generates a targeted diagnosis for an intermediate step, pointing out the error's location, type, and cause, without requiring an attached numeric score. DeepCritic~\cite{Yan25u} first does critique-teaching SFT on deliberate critiques generated by a teacher model, then does critique-incentivization GRPO on step-correctness data labeled automatically by PRM800K and MC, teaching the model to write multi-perspective, self-correcting long-form step annotations. AutoMathCritique~\cite{Xi24} has an actor generate flawed reasoning paths on GSM8K and MATH, has GPT-4o label the first error step, and keeps---through Monte Carlo soft filtering---the critiques that genuinely aid correction for SFT, yielding step-level feedback that can drive actor refinement directly.

\textbf{Verdict with structured reasoning} generates the supporting reasoning together with a verdict (correct/incorrect, promising/unpromising), making the verdict auditable and contestable. Unlike pure critique, this kind of evaluator must produce both verdict and reasoning under a logical constraint between them---the reasoning must support the verdict. StepWiser~\cite{Xio25c} has the model write meta-reasoning before giving a Positive/Negative verdict on a reasoning chunk, supervised by the chunk's before/after Q-value change estimated through MC continuation. OS-Oracle~\cite{Wu25m} synthesizes negatives by injecting four error types and uses CP-GRPO to reinforce the consistency between verdict and rationale. WebArbiter~\cite{Zha26ab} introduces a key change in the reasoning structure: before comparing candidate actions, the model must first induce principles applicable to the current task and page, then derive a preference verdict from those principles rather than jumping from state to judgment. Web-Shepherd~\cite{Cha25b} introduces another structure---checklist-grounded evaluation---generating a Yes/No/In-Progress judgment item by item against a predefined checklist, so the evaluation decomposes into individually verifiable sub-questions. SOLE-R1~\cite{Sch26} extends structured reasoning to the embodied temporal domain, learning from robot videos to generate per-timestep progress scores and a spatiotemporal CoT explaining its judgment. PPRM~\cite{Xu25h} generates step scores by principles such as correctness, relevance, and consistency on non-verifiable multi-turn search tasks, calibrated through ReNorm and outcome reward.

\textbf{Pre-operative corrective feedback} generates a diagnosis and corrective suggestion before an action executes, intercepting an error before it happens. GUI-Critic-R1~\cite{Wan25q} generates, in turn, an observation analysis, possible result, critique, correctness score, and corrective suggestion before a candidate action executes, reinforced by Suggestion-aware GRPO so that the critic's suggestions actually lower the subsequent execution-failure rate.

\subsubsection{Discussion}

Two supervision granularities and two product forms make four cells, and the literature is markedly uneven across them. Outcome labels are easy to obtain from environment signals, execution results, or human judges, and for tasks with a clear success criterion they align directly with the final goal; the outcome-discriminative cell therefore inherits the classic reward-modeling interface directly, is the most engineering-mature, and has the thickest literature---it is almost the default for episode-level success detection in embodied settings, and its downstream consumption lands directly on a scalar reward. The process-discriminative cell expands quickly around credit assignment, its axis of evolution being a semantic refinement from correctness to advantage, progress, Q-value, and bidirectional reward, then landing in domains with explicit intermediate states such as GUI, computer-use, code, and embodied; the cost is the acquisition cost (human annotation or MC estimation) and the noise of step-level labels. The two generative cells let the Evaluator output a readable critique or diagnosis, which downstream can collapse to a scalar for the RL of Evaluator-Driven Optimization, take as a binary for Best-of-N selection, or keep as full text to drive step-wise refinement---there is a monotone relation between product form and downstream utilization freedom, with the generative form offering the most. The cost is that the correctness of a critique lacks automatic verification---a plausible-sounding critique may be inconsistent with the real environment feedback---which is the open problem that Evaluator Internalization leaves to later work and cannot be resolved from the literature of this section alone.

Across supervision granularity, process supervision buys finer credit assignment at the price of annotation cost and noise; across product form, the generative form buys richer downstream utilization at a higher inference cost. The choice along both axes is set jointly by the downstream task's verifiability and supervision budget, and no choice is universally optimal.

\subsection{Policy Internalization: Tokenized-to-Policy Transformation}
\label{sec:policy-internalization}

Policy Internalization internalizes discretized agent experience into Policy parameters, fixing the sequential-decision capability in the weights. The source is the various Tokenized artifacts---natural-language trajectories, tool-call records, code-edit histories, GUI operation sequences, reasoning chains; once internalized, the model decides in later tasks without retrieving, concatenating, or replaying the original experience. The source experience is written into the Policy parameters under three objectives: imitation-based, environment-reward-based, and preference-based.

The inclusion boundary of this section is set by the provenance of the supervision signal that enters the policy update. When the signal is non-parametric---an environment-verifiable result, a unit-test pass, a ground-truth match, a rule-based verifier---the experience is transformed into Policy parameters directly, and the work belongs here. When the signal comes from a parametric evaluator (a trained RM/PRM or an LLM-as-a-judge), the policy update is mediated by the Evaluator rather than anchored to the experience directly: the experience is first internalized into an Evaluator and then drives the update through Evaluator-Driven Optimization, forming a P4 $\rightarrow$ P6 composite pipeline handled in Section~\ref{sec:composite}. Two cases still belong here: when a parametric evaluator is used only to filter or reconstruct the training data and its score does not enter the policy-update loss, that evaluator is not on the experience-transformation chain; and the value critic in actor-critic is bootstrapped internally by the environment reward within the algorithm, not an independently trained external reward source.

\subsubsection{Imitation-based Policy Internalization}

Imitation-based internalization treats high-quality agent trajectories as supervised demonstrations and writes the demonstrated action sequences into Policy parameters through behavior cloning, SFT, or rejection-sampled fine-tuning. The loss is always the maximum likelihood on successful data; environment feedback does not enter the optimization objective and serves only data selection and data construction---filtering out failed rollouts, removing invalid tool calls, keeping code edits that pass the tests. Subclass membership turns on the objective used to change the policy weights, regardless of whether the method is named RFT, ReST, or RL.

The main line is success-filtered behavior cloning. One approach has a web agent self-generate action--observation trajectories in WebArena, filters trustworthy actions by environment errors and self-critique, and does autoregressive SFT under QLoRA~\cite{Pat24}; SWE-Gym~\cite{Pan24} swaps the filter for executable unit tests, judging the success of a patch trajectory by test pass before rejection-sampling fine-tuning. Naming sometimes departs from the objective: AppVLM~\cite{Pap25b} calls itself RFT/ReST, but its final update is success-only MLE, with failed trajectories assigned return $0$ and excluded from the training set; the embodied VLA of \cite{Guo25d} calls itself online RL, but the RL stage only explores on the action head, and the final full-model update is supervised learning on the union of expert and newly discovered successful trajectories. When the filtering signal comes from a parametric evaluator but the evaluator only selects data and does not enter the loss, the work is still imitation: Go-Browse~\cite{Gan25b} gets its success labels from a VLM-as-a-judge and SFTs only successful trajectories, NNetNav~\cite{Mur24b} rewrites interactions through hindsight relabeling and then filters trajectories with an ORM, and AndroidGen~\cite{Lai25d} uses StepCritic to cut out successful subsegments before LoRA-SFT.

The demonstration source can be moved entirely to a stronger teacher. Explorer~\cite{Pah25} synthesizes multimodal web trajectories with a GPT-4o multi-agent pipeline and, after a Task Verifier filters them, has the student model learn the next action; AgentTuning~\cite{Zen23d} mixes multi-environment teacher trajectories with general instruction data under an NLL objective, so the model acquires the agent behavior protocol while keeping general instruction-following.

Trajectory processing before imitation falls into two memberships by whether it changes the carrier level. Light processing that reorders or relabels without changing the carrier---AndroidGen's subsegment cutting, NNetNav's hindsight relabeling---is preprocessing internal to imitation. A preprocessing step that reconstructs the experience across carriers is a different matter: WebCoT~\cite{Hu25i} verbalizes reflection, branching, and rollback traces into explicit CoT before SFT, and AgentTrek~\cite{Xu24} compiles a web tutorial into a Goal-Observation-Reasoning-Action schema before SFT---the former a Narrative $\rightarrow$ Narrative, the latter a Narrative $\rightarrow$ Schematic transformation, chained with imitation into a P1 $\rightarrow$ P5 or P2 $\rightarrow$ P5 composite pipeline whose contribution lies in the pathway junction, and which belongs to Section~\ref{sec:composite}.

\subsubsection{Environment-Reward-based Policy Internalization}

Environment-Reward-based internalization has the agent interact with the environment to produce trajectories, and feeds a non-parametric, programmatically verifiable signal directly into the policy-optimization objective. Verifiable signals cover results the environment itself can decide---task success/failure, unit-test pass/fail, execution output, SQL execution correctness, ground-truth answer match, GUI task-completion status, robot-manipulation success/failure; both successful and failed trajectories enter optimization, with the reward turning the former into a positive signal and the latter into a negative one. The difficulty arises in long-horizon multi-turn real environments: a trajectory of dozens of steps receives only a single sparse verifiable reward at the end, a signal hard to exploit---around which work proceeds in three directions: exploration preservation, attribution refinement, and reward densification.

Basic feasibility is established by two early works. GLAM~\cite{Car23} places FLAN-T5 in BabyAI-Text and grounds it online through PPO with a sparse goal reward; LLaRP~\cite{Szo23} connects a frozen LLaMA to a visual encoder and trains adapters and heads through DD-PPO with an environment reward composed of success, subgoal, and invalid-action penalties. When actions are executable and the environment returns a verifiable result, a Tokenized trajectory can be transformed into Policy parameters directly through RL.

\paragraph{Exploration preservation.} A long-horizon agent collapses easily to a low-exploration policy early in training, and once it collapses it stops producing successful trajectories and the gradient vanishes with them. EPO~\cite{Xu25k} adds trajectory-level entropy regularization and entropy smoothing outside PPO or GRPO, mitigating the exploration--exploitation cascade failure of sparse-reward multi-turn RL. Extensions along entropy control include Agentic Reinforced Policy Optimization~\cite{Don25d}, which triggers local branching at high-uncertainty tool steps; AEPO~\cite{Don25c}, which folds entropy into both advantage and clipping and allocates the global and branch rollout budget dynamically; and EMPG~\cite{Wan25ad}, which scales each step's advantage by the current entropy and adds a future-clarity bonus that rewards entering lower-entropy successor states. Another line keeps the success signal from drying up: ARPO~\cite{Lu25f} maintains a successful-trajectory replay buffer in GUI environments and injects positive samples back when a sampled group fails entirely and has no gradient; SoLS~\cite{Pap25c} updates strongly on positive-advantage samples and applies PPO-like conservative regularization to negative-advantage samples in mobile-app control, paired with successful-transition replay (STR); AgentGym-RL~\cite{Xi25c} uses ScalingInter-RL, a short-then-long horizon curriculum, so the agent first gets a signal on completable short-horizon tasks.

\paragraph{Fine-grained credit assignment.} The attribution granularity descends from the whole trajectory to the turn, step, and token. WebAgent-R1~\cite{Wei25} samples browser trajectories in parallel with M-GRPO and turns a binary success reward into a group-relative advantage; another work does token-level DAPO in a SWE environment with unit-test pass/fail plus a trajectory-length penalty~\cite{Gol25}. The actor-critic line uses a learned value to estimate intermediate advantage: ArCHer~\cite{Zho24f} first trains an utterance-level Q critic with the environment reward and then passes the multi-turn advantage to a token actor; Turn-PPO~\cite{Li25ae} models multi-turn interaction as a turn-level MDP and uses a critic for GAE at turn granularity. The critic-free line estimates relative advantage from group comparison: GiGPO~\cite{Fen25c} uses anchor-state grouping to compare the actions of different rollouts at the same intermediate state, building both episode-level and step-level relative advantages; LOOP~\cite{Che25af} does token-level updates with leave-one-out PPO in the AppWorld API environment. A further approach attaches the signal to intermediate semantics: IGPO~\cite{Wan25y} defines an information-gain intrinsic reward from the increment in the ground-truth answer's log-probability before and after each interaction turn, turning the real informational contribution of an observation into a turn-level dense signal; RLVMR~\cite{Zha25ao} tags trajectories with \texttt{<planning>}, \texttt{<explore>}, \texttt{<reflection>}, and \texttt{<monitor>} and uses a rule-verifiable process reward plus outcome reward to supervise meta-reasoning behavior explicitly.

\paragraph{Verifiable reward densification.} This direction reshapes the reward signal itself: mining a programmatically verifiable intermediate structure from the environment and writing a single terminal flag into a dense composite reward available at every step. ToolRL~\cite{Qia25b} splits the tool-use reward into format correctness and a three-level match of tool name, argument name, and argument value. The text-to-SQL domain is the most concentrated: SQL-Trail~\cite{Hua26d} composes a six-term composite reward---execution, turn budget, schema grounding, and others---within a reason-execute-observe loop; SQL-ASTRA~\cite{Li26r} gives dense partial credit from a Column-Set Matching Reward over result-table column-value sets and aggregates the whole trajectory with an Asymmetric Trajectory Reward. LongNav-R1~\cite{Hu26e} uses a geodesic-progress shaped reward plus a terminal SPL in long-horizon navigation, with a critic-free time-aware baseline for a REINFORCE++-style update. The GUI and VLA domains feed a rule-based verifiable reward straight into GRPO: GUI-R1~\cite{Luo25b}, AgentCPM-GUI~\cite{Zha25an}, and SimpleVLA-RL~\cite{Li25aa} stack decidable terms---action type, click point, text match, bbox hit---on top of format legality; UI-S1~\cite{Lu25j} uses model-self-generated history in semi-online RL, continuing with a patch module when it deviates from the expert, the reward provided by offline ground-truth matching. VLM as Decision-Making Agents~\cite{Zha24s} treats CoT plus action as an end-to-end text rollout and does PPO with an environment reward after $\lambda$-reweighting the log-probability contribution of the reasoning tokens. The same paradigm extends to multi-task and generalization studies: AgentRL~\cite{Zha25ag} stabilizes multi-task asynchronous RL with deterministic verifiable rewards plus task-advantage normalization, and Generalization in Mobile Agents~\cite{Gu26b} verifies instance-, template-, and app-level generalization on an asynchronous emulator. Q-SFT~\cite{Hon24} sits at the formal boundary of this subclass---formally close to SFT, but with the sample weight given by a Bellman target / Q-value, a mechanism closer to offline value-based RL, the reward provenance still an environment signal.

\subsubsection{Preference-based Policy Internalization}

Preference-based internalization builds a relative preference between trajectories or between steps from agent experience and updates the Policy directly with DPO or another trajectory-pair optimization. The source signal is a relative order: the advantaged item comes from a successful task, a verified execution, a corrected trajectory, or a more efficient operation path, and the disadvantaged item from a failed attempt, a wrong tool call, or a code edit that fails the tests. The granularity of the preference descends from the whole trajectory to local units within a trajectory.

\paragraph{Trajectory-level preference.} The preference is defined over a whole trajectory. DMPO~\cite{Shi24c} extends DPO from single-turn response preference to multi-turn trajectory preference, introducing a state-action occupancy measure and length normalization to handle the preference objective for long trajectories. ETO~\cite{Son24} first does behavior cloning on expert trajectories, then has the policy explore actively, pairing successful demonstrations with explored failed trajectories into failure--success pairs for continued DPO. Agent-RLVR~\cite{Da25} uses guidance in a SWE environment to help the model explore repair trajectories more likely to succeed, then builds preference pairs by unit-test pass/fail for DPO---here the guidance is a data generator and does not enter the reward.

\paragraph{Sub-trajectory preference.} The preference descends to local units within a trajectory. One method treats ``a successful step under the same parent node beats an erroneous branch step'' as a step-level preference pair in a tool-use inference tree, continuing with DPO after SFT~\cite{Che24k}. HPL~\cite{Gao25c} introduces preferences at three levels---trajectory, step, and action group, with groups carved out by fixed, entropy, or semantic segmentation---and its final objective is a composite of BC and three-level DPO, in which group-DPO is the main source of gain. WEPO~\cite{Liu25z} contracts the preference to the web-element level, training an element-selection policy by contrasting the target element against distractor elements close to it in DOM distance, with labels from the benchmark ground truth and DOM structure rather than a learned judge.

\subsubsection{Discussion}

The three classes differ in how a policy update uses a trajectory: take only its successful actions, take its full reward with signs, or take the relative order between trajectories. This choice simultaneously sets each one's signal-carrying capacity, training stability, and capability ceiling---the three are not independent knobs. Imitation-based does maximum likelihood on successful action sequences, and in the cold-start phase it teaches the action space, output protocol, and basic interaction flow fastest; because the supervision comes only from successful demonstrations, a strategy outside the demonstrated coverage cannot be acquired through MLE, and a limited set of success paths leaves the model in surface formats or local heuristics, with failed trajectories discarded wholesale. Environment-Reward-based has the reward carry the signs of both success and failure, so trial-and-error can discover strategies beyond the demonstrations and the capability ceiling is the highest; the same open-ended trial-and-error setting also means credit assignment is hard when the reward is sparse, online rollout is expensive, reward hacking is hard to prevent, and once the high-variance RL gradient collapses or is hacked, the damaged capability cannot be rolled back item by item the way a textual rule can. Preference-based takes the relative order between trajectories: a relative order is a supervised objective, needs no high-variance policy-gradient estimate, trains more stably than RL, and brings failed trajectories in as negatives rather than discarding them; the cost is that a relative order only constrains the ranking of a pair and gives no absolute signal, so the effect rests on the construction quality of the preference pairs, and a trajectory-level preference may still lack fine-grained attribution.

The three are often chained as different stages of one pipeline: imitation stands up the format and protocol, and reward optimization or preference learning takes over the capability breakthrough---they are therefore more complementary than mutually exclusive (the integration analysis of such stage chaining is in Section~\ref{sec:composite}). Environment-Reward and Preference descend along the same credit-assignment axis---the former toward turn/step/token, the latter toward step/element/group---and divide on the supervision form: an absolute verifiable reward enters the RL objective, a relative pair enters the supervised loss. The three also share a property of the Parametric carrier: once experience is written into the weights it cannot be addressed locally, and adding or correcting a class of behavior requires retraining rather than the entry-by-entry editing a Tokenized carrier allows.
